% from ega4-16-fr.tex
\oldpage[IV-4]{39}

\begin{env}[16.7.8]
\label{IV.16.7.8}
Soient $\mathscr{F}$, $\mathscr{G}$ deux $\mathcal{O}_X$-Modules; il résulte
aussitôt de la définition (16.7.2.1) que l'on a un isomorphisme canonique de
$\mathscr{P}_{X/S}^n$-Modules
\[
  \mathscr{P}_{X/S}^n(\mathscr{F}\otimes_{\mathcal{O}_X}\mathscr{G})
  \xrightarrow{\sim}
  \mathscr{P}_{X/S}^n(\mathscr{F})
  \otimes_{\mathscr{P}_{X/S}^n}\mathscr{P}_{X/S}^n(\mathscr{G})
  \tag{16.7.8.1}
  \label{IV.16.7.8.1}
\]
(Bourbaki, \emph{Alg.}, chap.~II, 3\textsuperscript{e} éd., §~5,
n\textsuperscript{o}~1, prop.~3).

On en conclut en particulier (ou l'on voit directement sur la définition
(16.7.2.1)) que si $\mathscr{F}$ est muni d'une structure de
$\mathcal{O}_X$-Algèbre (non nécessairement associative),
$\mathscr{P}_{X/S}^n(\mathscr{F})$ est canoniquement muni d'une structure de
$\mathscr{P}_{X/S}^n$-Algèbre; cette dernière est associative (resp.
commutative, resp. unitaire, resp. une Algèbre de Lie) lorsqu'il en est ainsi
de $\mathscr{F}$. En outre les homomorphismes canoniques
$\mathscr{P}_{X/S}^m(\mathscr{F})\to\mathscr{P}_{X/S}^n(\mathscr{F})$ pour
$n\leq m$ (16.7.7) sont alors des di-homomorphismes d'Algèbres; de même
(16.7.5.1) est alors un homomorphisme de $\mathcal{O}_X$-Algèbres lorsque
$\mathscr{P}_{X/S}^n(\mathscr{F})$ est muni de sa structure de
$\mathcal{O}_X$-Algèbre provenant de sa structure de $\mathcal{O}_X$-Module
\emph{à droite}.

Avec les mêmes notations, on a également un homomorphisme canonique de
$\mathscr{P}_{X/S}^n$-Modules
\[
  \mathscr{P}_{X/S}^n
  (\mathscr{H}\!om_{\mathcal{O}_X}(\mathscr{F},\mathscr{G}))
  \longrightarrow
  \mathscr{H}\!om_{\mathscr{P}_{X/S}^n}
  (\mathscr{P}_{X/S}^n(\mathscr{F}),\mathscr{P}_{X/S}^n(\mathscr{G}))
  \tag{16.7.8.2}
  \label{IV.16.7.8.2}
\]
(Bourbaki, \emph{Alg.}, 3\textsuperscript{e} éd., §~5,
n\textsuperscript{o}~3), qui est bijectif lorsque $\mathscr{P}_{X/S}^n$ est
un $\mathcal{O}_X$-Module \emph{localement libre de type fini}
(\emph{loc. cit.}, prop.~7).
\end{env}

\begin{env}[16.7.9]
\label{IV.16.7.9}
Supposons qu'on soit dans la situation décrite dans (16.4.1); alors, de
l'homomorphisme canonique $P^n(u)$ (16.4.3.3), on déduit aussitôt un
homomorphisme canonique de $\mathcal{O}_{X'}$-Bimodules
\[
  u^*(\mathscr{P}_{X/S}^n(\mathscr{F}))
  \longrightarrow
  \mathscr{P}_{X'/S'}^n(u^*(\mathscr{F})).
  \tag{16.7.9.1}
  \label{IV.16.7.9.1}
\]
Nous laissons au lecteur le soin d'étendre à cet homomorphisme les propriétés
vues dans (16.4) pour le cas $\mathscr{F}=\mathcal{O}_X$.
\end{env}

\begin{remark}[16.7.10]
\label{IV.16.7.10}
La définition de $\mathscr{P}_{X/S}^n(\mathscr{F})$ sous la forme (16.7.1.2)
garde un sens lorsque $\mathscr{F}$ est un faisceau d'ensembles quelconque
(l'image réciproque d'un faisceau d'ensembles par $p_2^{(n)}$ étant définie
dans (0\textsubscript{I}, 3.7.1)); une variante de cette définition permet de
définir le «schéma des jets» (relativement à $S$) d'un $X$-préschéma
quelconque.
\end{remark}

\subsection{Opérateurs différentiels.}
\footnote{Pour un formalisme plus général, voir l'Exposé VII de [42]
(dû à P. Gabriel).}
\EGAReaderTarget{ega4.c04.subsec.16-8}{IV.16.8}

\begin{definition}[16.8.1]
\EGAReaderTarget{ega4.c04.def.16-8-1}{IV.16.8.1}
Soient $f=(\psi,\theta):X\to S$ un morphisme de préschémas,
$\mathscr{F}$, $\mathscr{G}$ deux $\mathcal{O}_X$-Modules, $n$ un entier
$\geq0$. On dit qu'un homomorphisme de faisceaux de groupes additifs
$D:\mathscr{F}\to\mathscr{G}$ est un \emph{opérateur différentiel d'ordre
$\leq n$} (relativement à $S$) s'il existe un homomorphisme de
$\mathcal{O}_X$-Modules
$u:\mathscr{P}_{X/S}^n(\mathscr{F})\to\mathscr{G}$ (où
$\mathscr{P}_{X/S}^n(\mathscr{F})$ est muni de sa structure de
$\mathcal{O}_X$-Module à gauche (16.7.4)) tel que l'on ait
$D=u\circ d_{X/S,\mathscr{F}}^n$.
\end{definition}

Il est clair, en vertu de l'existence des homomorphismes canoniques
\[
  \mathscr{P}_{X/S}^m(\mathscr{F})
  \longrightarrow\mathscr{P}_{X/S}^n(\mathscr{F})
\]
